\documentclass[times]{elsarticle}
% \documentclass[review,doubleblind,times,12p]{elsarticle}
\let\today\relax

\usepackage[margin=2.5cm]{geometry}
\usepackage{amsmath,amssymb,bm}
\usepackage{algorithm}
\usepackage[noend]{algpseudocode}
\usepackage{booktabs,enumitem}
\usepackage{xurl}
\usepackage[colorlinks=true,citecolor=blue,linkcolor=blue,urlcolor=blue]{hyperref}
\makeatletter
\pdfstringdefDisableCommands{%
  \let\corref\@gobble
  \let\cnotenum\@empty
  \let\@corref\@empty
}
\makeatother

\allowdisplaybreaks
\setlength{\parindent}{12pt}
\DeclareMathOperator*{\argmin}{\arg\!\min}

\newenvironment{manuscriptdescription}
  {\begin{description}[style=multiline,leftmargin=4.8cm,labelwidth=4.5cm,align=left]}
  {\end{description}}

\begin{document}
\begin{frontmatter}

\title{Feasibility-Aware Scenario Learning for Carbon-Constrained Battery Planning in Data Centers}

\author[seuaddress]{Author Name\corref{correspondingauthor}}
\cortext[correspondingauthor]{Corresponding author.}
\ead{author@seu.edu.cn}
\address[seuaddress]{School of Electrical Engineering, Southeast University,
Nanjing, China}

\begin{abstract}
Planning battery storage for a data center requires a single design to balance
annual electricity cost, outage service, and the carbon intensity of delivered
electricity. The design can be determined by rare combinations of high demand,
low photovoltaic availability, carbon-intensive grid supply, and grid outages;
scenario supports selected for statistical similarity need not preserve these
decision-critical events. This paper develops a two-stage battery-planning model
and a feasibility-aware method for learning a compact scenario support. The
planning model uses a source- and charging-time-resolved storage ledger that
tracks initial, photovoltaic-, diesel-, and grid-charged energy without a
bilinear perfect-mixing approximation. A conditional variational autoencoder is
then fine-tuned using fixed-design recourse from the mixed-integer planning
oracle. The decision loss combines an infeasibility-penalizing term, which
corrects optimistic demand, photovoltaic, and carbon trajectories, with an
optimality-preserving term, which prevents uniformly conservative scenarios
from excluding the design preferred for matched real data. Solver outputs
enter the loss only as detached violation severities and carbon-exposure
weights; neither the optimizer nor decision regret is differentiated. The
evaluation compares learned, statistical-selection, and predict-then-optimize
supports using common held-out scenarios, installed capacity, annual cost, load
shedding, carbon excess, and decision regret.
% TODO: Replace the final sentence with quantitative results.
\end{abstract}

\begin{keyword}
Battery energy storage \sep carbon accounting \sep data center \sep
decision-focused learning \sep scenario generation \sep stochastic planning
\end{keyword}

\end{frontmatter}

% =====================================================================
\section*{Nomenclature}

\subsubsection*{Sets and Indices}
\begin{manuscriptdescription}
\item[$\mathcal S$] Planning scenarios, indexed by $s$.
\item[$\mathcal T,\mathcal T^+$] Operating intervals and storage-state intervals,
indexed by $t$.
\item[$\mathcal J$] Supply sources $\{\mathrm{G},\mathrm{PV},\mathrm{DG}\}$,
indexed by $j$.
\item[$\mathcal K$] Source- and vintage-resolved storage layers, indexed by $k$.
\item[$\mathcal I$] Generated-support positions, indexed by $n$.
\item[$\mathcal D^{\mathrm{tr}}$] Historical training pool.
\item[$\mathrm L,\mathrm{ES}$] Facility-load and energy-storage destinations.
\end{manuscriptdescription}

\subsubsection*{Scenario and System Parameters}
\begin{manuscriptdescription}
\item[$D_{s,t},P^{\mathrm{PV,av}}_{s,t}$] Facility demand and available PV power.
\item[$\lambda_{s,t},\chi_{s,t}$] Grid energy price and carbon intensity.
\item[$A_{s,t}$] Grid-availability indicator.
\item[$\pi_s,\nu_s$] Physical probability and annual occurrence factor.
\item[$\Delta t$] Duration of one operating interval.
\item[$\eta^{\mathrm{ch}},\eta^{\mathrm{dc}}$] Storage charging/discharging efficiencies.
\item[$\sigma^{\mathrm{sd}},\kappa$] Self-discharge fraction and interval retention
factor.
\item[$\sigma^0,\underline\sigma,\overline\sigma$] Initial, minimum, and maximum
state-of-charge fractions.
\item[$\underline P^{\mathrm{ES}},\overline P^{\mathrm{ES}}$] Minimum/maximum
storage power if built.
\item[$\underline E^{\mathrm{ES}},\overline E^{\mathrm{ES}}$] Minimum/maximum
storage energy if built.
\item[$\underline h,\overline h$] Minimum/maximum storage duration.
\item[$\overline P^{\mathrm G},\overline P^{\mathrm{DG}}$] Grid-import and
diesel-generator limits.
\item[$\chi^0,\chi^{\mathrm{DG}}$] Initial-inventory and diesel carbon intensities.
\item[$\bar\chi^{\mathrm n},\bar\chi^{\mathrm o}$] Normal-operation and outage
carbon-intensity limits.
\item[$\chi^{\mathrm{ES}}_{s,k}$] Internal carbon intensity of storage layer $k$.
\end{manuscriptdescription}

\subsubsection*{Cost Parameters}
\begin{manuscriptdescription}
\item[$\mathrm{CRF}$] Capital recovery factor.
\item[$c^{\mathrm{site}},c^P,c^E$] Fixed, power-related, and energy-related
storage investment costs.
\item[$c^{\mathrm{dem}},c^{\mathrm{DG}}$] Annual demand-charge rate and diesel
energy cost.
\item[$c^{\mathrm{deg}},c^{\mathrm{cur}},c^{\mathrm{sh}}$] Storage degradation,
PV-curtailment, and load-shedding costs.
\item[$c^{\mathrm{exc}}$] Penalty per unit of carbon-budget excess.
\end{manuscriptdescription}

\subsubsection*{Planning and Operating Variables}
\begin{manuscriptdescription}
\item[$u^{\mathrm{ES}},P^{\mathrm{ES}},E^{\mathrm{ES}}$] Storage build,
power-capacity, and energy-capacity decisions.
\item[$p^{\mathrm{pk}}$] Maximum grid import over the planning support.
\item[$p^{j,\mathrm L}_{s,t},p^{j,\mathrm{ES}}_{s,t}$] Power from source $j$
to the load and storage.
\item[$p^{\mathrm{ch}}_{s,t},p^{\mathrm{dc}}_{s,t}$] Aggregate storage charging and
discharging power.
\item[$p^{\mathrm{sh}}_{s,t},p^{\mathrm{cur}}_{s,t}$] Load shedding and PV curtailment.
\item[$u^{\mathrm{ch}}_{s,t},u^{\mathrm{dc}}_{s,t}$] Storage charge/discharge
mode indicators.
\item[$e_{s,k,t},p^{\mathrm{dc}}_{s,k,t}$] Energy and discharge power of layer $k$.
\item[$\Delta e^{\mathrm{in}}_{s,k,t}$] Internal energy added to layer $k$.
\item[$m^{\mathrm L}_{s,t},\varepsilon^{\mathrm{co2}}_{s,t}$] Carbon delivered to load and hourly carbon
excess.
\end{manuscriptdescription}

\subsubsection*{Learning Symbols}
\begin{manuscriptdescription}
\item[$G_\theta$] Conditional scenario decoder.
\item[$\widehat{\mathcal S}_\theta$] Scenario support generated by $G_\theta$.
\item[$\widehat{\bm y}_\theta,\bm y^{\star}$] Generated-support and matched-real
reference designs.
\item[$\delta^{\mathrm{sh}},\delta^{\mathrm{co2}}$] Detached shedding and carbon-excess
severities returned by recourse.
\item[$H_{n,t,\tau}$] Detached exposure of delivered carbon in $t$ to grid
carbon intensity in $\tau$.
\item[$F,P$] Forward-feasibility and preservation-recourse superscripts.
\item[$\epsilon_D,\epsilon_H,\epsilon_g,\epsilon_J$] Positive numerical safeguards.
\item[$L_{\mathrm{IPL}},L_{\mathrm{OPL}}$] Infeasibility-penalizing and
optimality-preserving losses.
\item[$\mathcal R_m$] Normalized out-of-sample decision regret of method $m$.
\end{manuscriptdescription}

% =====================================================================
\section{Introduction}
\label{sec:introduction}

Data centers combine large, persistent electricity demand with stringent
service-continuity requirements. Behind-the-meter photovoltaic (PV) generation
and battery storage can reduce demand charges, shift grid imports toward
lower-carbon hours, and sustain critical demand during an outage. These services
cannot be planned independently. A battery selected for price arbitrage may
lack sufficient power or energy during an interruption, whereas a design chosen
only for reliability may be charged in carbon-intensive hours. The planning
problem must therefore represent economic performance, outage operation, and
delivered carbon under the same uncertain trajectories.

Carbon-aware storage creates an accounting problem in addition to the physical
energy balance. Grid electricity charged at different times can have different
carbon intensities, and PV-, diesel-, and grid-charged electricity should not
lose its identity merely because it occupies one physical battery. A perfectly
mixed inventory assigns a single endogenous average intensity to all stored
energy. Besides introducing bilinear relations, that abstraction discards the
source and charging-time information that determines the carbon content of a
later discharge. Research on carbon-aware power flow has shown why emissions
must be attributed to delivered electricity~\cite{chen2025copf}; storage extends
the same attribution question across time.

Uncertainty creates a second difficulty. A planning scenario must jointly
represent facility demand, PV availability, energy price, grid carbon
intensity, and grid availability. Historical data provide many such
trajectories, but repeatedly solving the resulting mixed-integer planning model
on the complete pool is computationally burdensome. Classical scenario
reduction approximates the empirical distribution under a statistical distance
~\cite{Dupacova2003}. A statistically representative support, however, need not
preserve the events that determine storage capacity. An error concentrated in
an outage hour or in a carbon-intensive charging hour can change the design even
when its contribution to average reconstruction error is small.

Decision-focused learning aligns prediction with downstream optimization
~\cite{Wilder2019}, but the present problem does not admit a convenient
end-to-end derivative. The planning oracle contains build and operating
binaries, and only the first-stage battery design transfers from generated
scenarios to realized operation. Moreover, generated demand, PV availability,
and grid carbon intensity appear in feasibility constraints. An optimistic
support can produce an inexpensive design that sheds load or exceeds its carbon
budget under real recourse; a uniformly conservative support can avoid these
violations while excluding an otherwise effective design. This distinction
motivates the feasibility/optimality-preservation framework in
~\cite{Mandi2025}.

This paper combines exact storage planning with conditional scenario
generation. A conditional variational autoencoder (CVAE)~\cite{Sohn2015}
produces a fixed support of 48-hour trajectories. The planning oracle selects a
battery from that support, after which the design is fixed and every operating
variable is reoptimized on matched real trajectories. Detached load-shedding,
carbon-excess, and carbon-exposure outputs identify optimistic generated
constraints. In the reverse direction, a design optimized for the matched real
batch is fixed and re-dispatched on the generated support to identify excessive
conservatism.

The main contributions are as follows.
\begin{enumerate}[leftmargin=*,label=\arabic*.]
\item A two-stage data-center battery-planning model jointly represents
investment, energy and demand charges, outage-only load shedding, and hourly
load-side carbon limits. A linear source- and charging-time-resolved ledger
preserves the carbon intensity of stored electricity without the bilinear
identity required by a perfectly mixed inventory.
\item A feasibility-aware scenario-learning method is developed for
mixed-integer, two-stage planning. Fixed-design recourse generates
infeasibility-penalizing and optimality-preserving signals, while detached
solver outputs avoid differentiation through the planning oracle.
\item A controlled evaluation protocol uses disjoint data splits, identical
support sizes, common held-out recourse, and a perfect-information reference to
compare annual cost, physical violations, installed capacity, storage value,
and decision regret.
\end{enumerate}

Section~\ref{sec:formulation} formulates the storage-planning problem.
Section~\ref{sec:methodology} presents the decision-focused methodology.
Section~\ref{sec:case} describes the case study, and
Section~\ref{sec:conclusion} concludes the paper.

% =====================================================================
\section{Carbon-Aware Storage Planning Formulation}
\label{sec:formulation}

\subsection{Problem Definition and System Components}

Consider a \emph{carbon-aware storage planning} (CASP) problem for a
grid-connected data center subject to an operational carbon-intensity limit.
The CASP is formulated as a two-stage stochastic program. The first stage
optimizes the battery investment, i.e., whether to install a battery and its
power and energy ratings. For each scenario
$s\in\mathcal S=\{1,\ldots,S\}$, the second stage optimizes system operation
over the operating horizon $\mathcal T=\{1,\ldots,T\}$, subject to load-service
and carbon-intensity constraints.

The system comprises an aggregate data-center load ($\mathrm L$), the utility
grid ($\mathrm G$), behind-the-meter photovoltaic generation ($\mathrm{PV}$),
a dispatchable diesel generator ($\mathrm{DG}$), and a candidate battery energy
storage system ($\mathrm{ES}$). The supply sources are indexed by
$\mathcal J=\{\mathrm G,\mathrm{PV},\mathrm{DG}\}$ and can serve the facility
load directly or charge the storage system. Storage states are indexed by
$\mathcal T^+=\{1,\ldots,T+1\}$.

\subsection{Planning Objective}

The first-stage design is represented by
\begin{equation}
\bm y=(u^{\mathrm{ES}},P^{\mathrm{ES}},E^{\mathrm{ES}}),
\label{eq:design_vector}
\end{equation}
where $u^{\mathrm{ES}}$ indicates whether storage is installed and
$P^{\mathrm{ES}}$ and $E^{\mathrm{ES}}$ are its rated power and energy. Its
annualized investment cost is
\begin{equation}
C^{\mathrm{inv}}(\bm y)=\mathrm{CRF}
\left(c^{\mathrm{site}}u^{\mathrm{ES}}+c^P P^{\mathrm{ES}}
+c^E E^{\mathrm{ES}}\right).
\label{eq:investment}
\end{equation}

Scenario $s\in\mathcal S$ is the joint trajectory
\begin{equation}
\bm\xi_s=
\{D_{s,t},P^{\mathrm{PV,av}}_{s,t},\lambda_{s,t},\chi_{s,t},A_{s,t}\}_{t\in\mathcal T},
\label{eq:scenario}
\end{equation}
where $D_{s,t}$ is demand, $P^{\mathrm{PV,av}}_{s,t}$ is available PV power,
$\lambda_{s,t}$ is the grid energy price, $\chi_{s,t}$ is the grid carbon
intensity, and $A_{s,t}\in\{0,1\}$ is the grid-availability indicator. If
$\pi_s$ is the physical probability of scenario $s$, its annual occurrence
factor is
\begin{equation}
\nu_s=\pi_s\frac{8760}{T\Delta t},
\label{eq:annualization}
\end{equation}
where $\pi_s\ge0$ and $\sum_{s\in\mathcal S}\pi_s=1$. Importance-corrected
probabilities prevent oversampled outage scenarios from inflating their assumed
annual frequency.

Let $p^{\mathrm{pk}}$ be the maximum total grid import over all scenarios. The
annual operating cost is
\begin{align}
C^{\mathrm{op}}={}&c^{\mathrm{dem}}p^{\mathrm{pk}}
+\sum_{s\in\mathcal S}\nu_s\Delta t\sum_{t\in\mathcal T}\big[
\lambda_{s,t}(p^{\mathrm{G},\mathrm L}_{s,t}
+p^{\mathrm{G},\mathrm{ES}}_{s,t})
\nonumber\\
&+c^{\mathrm{DG}}(p^{\mathrm{DG},\mathrm L}_{s,t}
+p^{\mathrm{DG},\mathrm{ES}}_{s,t})
+c^{\mathrm{deg}}(p^{\mathrm{ch}}_{s,t}+p^{\mathrm{dc}}_{s,t})
\nonumber\\
&+c^{\mathrm{cur}}p^{\mathrm{cur}}_{s,t}
+c^{\mathrm{sh}}p^{\mathrm{sh}}_{s,t}\big],
\label{eq:operating_cost}
\end{align}
Carbon excess, when enabled during training and diagnosis, incurs
\begin{equation}
C^{\mathrm{exc}}=
c^{\mathrm{exc}}\sum_{s\in\mathcal S}\nu_s\Delta t
\sum_{t\in\mathcal T}\varepsilon^{\mathrm{co2}}_{s,t}.
\label{eq:carbon_excess_cost}
\end{equation}
For the horizon-average benchmark, this term becomes
$C^{\mathrm{exc}}=c^{\mathrm{exc}}
\sum_{s\in\mathcal S}\nu_s\varepsilon^{\mathrm{co2}}_s$.
The planning oracle minimizes
\begin{equation}
C^{\mathrm{ann}}=C^{\mathrm{inv}}+C^{\mathrm{op}}+C^{\mathrm{exc}}.
\label{eq:annual_cost}
\end{equation}
When a hard carbon limit is required,
$\varepsilon^{\mathrm{co2}}_{s,t}$ is fixed to zero and $C^{\mathrm{exc}}$
vanishes. Soft-cap solves retain this variable so carbon infeasibility can be
penalized and compared on a common recourse problem.

\subsection{Storage Design and Site Dispatch}

The build decision links the installed capacities to their admissible ranges:
\begin{align}
\underline P^{\mathrm{ES}}u^{\mathrm{ES}}&\le P^{\mathrm{ES}}
\le\overline P^{\mathrm{ES}}u^{\mathrm{ES}},
\label{eq:power_cap}\\
\underline E^{\mathrm{ES}}u^{\mathrm{ES}}&\le E^{\mathrm{ES}}
\le\overline E^{\mathrm{ES}}u^{\mathrm{ES}},
\label{eq:energy_cap}\\
\underline h P^{\mathrm{ES}}&\le E^{\mathrm{ES}}
\le\overline h P^{\mathrm{ES}}.
\label{eq:duration}
\end{align}
Here, $u^{\mathrm{ES}}\in\{0,1\}$ and $p^{\mathrm{pk}}\ge0$. Thus,
$u^{\mathrm{ES}}$ determines
installation rather than location.

For source $j\in\mathcal J$, $p^{j,\mathrm L}_{s,t}$ supplies the facility
directly and $p^{j,\mathrm{ES}}_{s,t}$ charges storage. The facility and PV
balances are
\begin{align}
\sum_{j\in\mathcal J}p^{j,\mathrm L}_{s,t}+p^{\mathrm{dc}}_{s,t}
&=D_{s,t}-p^{\mathrm{sh}}_{s,t},
\label{eq:site_balance}\\
p^{\mathrm{PV},\mathrm L}_{s,t}+p^{\mathrm{PV},\mathrm{ES}}_{s,t}
+p^{\mathrm{cur}}_{s,t}
&=P^{\mathrm{PV,av}}_{s,t}.
\label{eq:pv_balance}
\end{align}
Grid availability, generator rating, and the demand-charge epigraph require
\begin{align}
p^{\mathrm{G},\mathrm L}_{s,t}+p^{\mathrm{G},\mathrm{ES}}_{s,t}
&\le A_{s,t}\overline P^{\mathrm G},
\label{eq:grid_limit}\\
p^{\mathrm{DG},\mathrm L}_{s,t}+p^{\mathrm{DG},\mathrm{ES}}_{s,t}
&\le\overline P^{\mathrm{DG}},
\label{eq:diesel_limit}\\
p^{\mathrm{G},\mathrm L}_{s,t}+p^{\mathrm{G},\mathrm{ES}}_{s,t}
&\le p^{\mathrm{pk}}.
\label{eq:peak_limit}
\end{align}
Load shedding is restricted to outages:
\begin{equation}
0\le p^{\mathrm{sh}}_{s,t}\le D_{s,t}(1-A_{s,t}).
\label{eq:outage_shedding}
\end{equation}
Consequently, a carbon limit cannot be met by dropping grid-connected load.

Aggregate charging and discharging are
\begin{align}
p^{\mathrm{ch}}_{s,t}&=\sum_{j\in\mathcal J}p^{j,\mathrm{ES}}_{s,t},
\label{eq:total_charge}\\
p^{\mathrm{dc}}_{s,t}&=\sum_{k\in\mathcal K}p^{\mathrm{dc}}_{s,k,t}.
\label{eq:total_discharge}
\end{align}
Binary variables $u^{\mathrm{ch}}_{s,t}$ and $u^{\mathrm{dc}}_{s,t}$ prevent
simultaneous charging and discharging:
\begin{align}
p^{\mathrm{ch}}_{s,t}&\le P^{\mathrm{ES}}, &
p^{\mathrm{ch}}_{s,t}&\le\overline P^{\mathrm{ES}} u^{\mathrm{ch}}_{s,t},
\label{eq:charge_limit}\\
p^{\mathrm{dc}}_{s,t}&\le P^{\mathrm{ES}}, &
p^{\mathrm{dc}}_{s,t}&\le\overline P^{\mathrm{ES}} u^{\mathrm{dc}}_{s,t},
\label{eq:discharge_limit}\\
u^{\mathrm{ch}}_{s,t}+u^{\mathrm{dc}}_{s,t}&\le u^{\mathrm{ES}}.
\label{eq:mode}
\end{align}
The mode variables are binary, and all power-flow variables are nonnegative.

\subsection{Source- and Vintage-Resolved Carbon Accounting}

To retain the origin of stored electricity, the battery inventory is partitioned
into the virtual layers
\begin{equation}
\mathcal K=\{0,\mathrm{PV},\mathrm{DG}\}
\cup\{\mathrm{G}_\tau:\tau\in\mathcal T\},
\label{eq:layers}
\end{equation}
where $0$, $\mathrm{PV}$, and $\mathrm{DG}$ denote initial, PV-charged, and
diesel-charged inventory, while $\mathrm{G}_\tau$ denotes grid energy charged
in interval $\tau$.
These layers share the same physical power and energy ratings; they are carbon
bookkeeping states, not separate battery compartments.

The internal carbon intensity assigned to each layer is
\begin{equation}
\chi^{\mathrm{ES}}_{s,k}=
\begin{cases}
\chi^0,&k=0,\\
0,&k=\mathrm{PV},\\
\chi^{\mathrm{DG}}/\eta^{\mathrm{ch}},&k=\mathrm{DG},\\
\chi_{s,\tau}/\eta^{\mathrm{ch}},&k=\mathrm{G}_\tau.
\end{cases}
\label{eq:layer_intensity}
\end{equation}
The division by $\eta^{\mathrm{ch}}$ expresses source emissions per unit of
internal stored energy: charging losses reduce energy entering the cell but do
not remove emissions already attributed to that electricity.

The energy added to layer $k$ in interval $t$ is
\begin{equation}
\Delta e^{\mathrm{in}}_{s,k,t}=
\begin{cases}
\eta^{\mathrm{ch}}p^{\mathrm{PV},\mathrm{ES}}_{s,t}\Delta t,
&k=\mathrm{PV},\\
\eta^{\mathrm{ch}}p^{\mathrm{DG},\mathrm{ES}}_{s,t}\Delta t,
&k=\mathrm{DG},\\
\eta^{\mathrm{ch}}p^{\mathrm{G},\mathrm{ES}}_{s,t}\Delta t,
&k=\mathrm{G}_t,\\
0,&\text{otherwise}.
\end{cases}
\label{eq:additions}
\end{equation}
Let $\kappa=1-\sigma^{\mathrm{sd}}$ be the one-interval retention factor. For
$s\in\mathcal S$, $k\in\mathcal K$, and $t\in\mathcal T$, layer energy evolves
according to
\begin{align}
e_{s,k,t+1}&=\kappa e_{s,k,t}+\Delta e^{\mathrm{in}}_{s,k,t}
-\frac{p^{\mathrm{dc}}_{s,k,t}}{\eta^{\mathrm{dc}}}\Delta t,
\label{eq:layer_energy}\\
\frac{p^{\mathrm{dc}}_{s,k,t}}{\eta^{\mathrm{dc}}}\Delta t
&\le\kappa e_{s,k,t}.
\label{eq:available_energy}
\end{align}
The initial layer allocation and aggregate state-of-charge bounds are
\begin{align}
e_{s,0,1}&=\sigma^0E^{\mathrm{ES}},
&e_{s,k,1}&=0,\quad k\ne0,
\label{eq:initial_layers}\\
\underline\sigma E^{\mathrm{ES}}
&\le\sum_{k\in\mathcal K}e_{s,k,t}\le\overline\sigma E^{\mathrm{ES}},
\quad t\in\mathcal T^+.
\label{eq:soc}
\end{align}
Each representative block returns to its initial stored energy and cannot end
with a larger carbon inventory:
\begin{align}
\sum_{k\in\mathcal K}e_{s,k,T+1}&=\sigma^0E^{\mathrm{ES}},
\label{eq:terminal_energy}\\
\sum_{k\in\mathcal K}\chi^{\mathrm{ES}}_{s,k}e_{s,k,T+1}
&\le \chi^0\sigma^0E^{\mathrm{ES}}.
\label{eq:terminal_carbon}
\end{align}

The carbon-delivery rate assigned to served facility load is
\begin{equation}
m^{\mathrm L}_{s,t}=\chi_{s,t}p^{\mathrm{G},\mathrm L}_{s,t}
+\chi^{\mathrm{DG}}p^{\mathrm{DG},\mathrm L}_{s,t}
+\sum_{k\in\mathcal K}\chi^{\mathrm{ES}}_{s,k}
\frac{p^{\mathrm{dc}}_{s,k,t}}{\eta^{\mathrm{dc}}}.
\label{eq:delivered_carbon}
\end{equation}
Carbon associated with charging enters the inventory ledger and is attributed
to the facility only when the stored energy is discharged. Hence,
$m^{\mathrm L}_{s,t}$ is a load-side compliance rate; it is not the source-side emissions
rate, which also counts grid and diesel energy used for charging.

\subsection{Carbon Limits and Benchmark Accounting}

The carbon-intensity limit depends on grid availability:
\begin{equation}
\bar\chi_{s,t}=A_{s,t}\bar\chi^{\mathrm n}
+(1-A_{s,t})\bar\chi^{\mathrm o}.
\label{eq:applicable_cap}
\end{equation}
The principal formulation imposes the following hourly load-side budget:
\begin{equation}
m^{\mathrm L}_{s,t}\le
\bar\chi_{s,t}(D_{s,t}-p^{\mathrm{sh}}_{s,t})
+\varepsilon^{\mathrm{co2}}_{s,t},
\qquad \varepsilon^{\mathrm{co2}}_{s,t}\ge0.
\label{eq:hourly_cap}
\end{equation}
For comparison, a horizon-average cap replaces (\ref{eq:hourly_cap}) with
\begin{equation}
\sum_{t\in\mathcal T}m^{\mathrm L}_{s,t}\Delta t\le
\sum_{t\in\mathcal T}\bar\chi_{s,t}
(D_{s,t}-p^{\mathrm{sh}}_{s,t})\Delta t+\varepsilon^{\mathrm{co2}}_s.
\label{eq:horizon_cap}
\end{equation}

The layered formulation is linear apart from the build and operating-mode
binaries. A perfectly mixed benchmark instead tracks energy
$E^{\mathrm{mix}}_{s,t}$, carbon inventory $M^{\mathrm{mix}}_{s,t}$, and
average internal intensity $\chi^{\mathrm{mix}}_{s,t}$ through
\begin{equation}
M^{\mathrm{mix}}_{s,t}=\chi^{\mathrm{mix}}_{s,t}E^{\mathrm{mix}}_{s,t},
\qquad
m^{B,\mathrm{dc}}_{s,t}=\chi^{\mathrm{mix}}_{s,t}
\frac{p^{\mathrm{dc}}_{s,t}}{\eta^{\mathrm{dc}}}.
\label{eq:average_model}
\end{equation}
These bilinear identities make the mixed-inventory benchmark nonconvex. The
benchmark is used only to determine whether discarding source and charging-time
information changes the resulting storage design.

\subsection{Compact Formulation}

Let $\bm x_s$ collect the recourse variables for scenario $s$, let
$\mathcal Y$ denote the design set in
(\ref{eq:power_cap})--(\ref{eq:duration}), and let
$\mathcal X(\bm y,\bm\xi_s)$ denote the feasible operating set defined by
(\ref{eq:site_balance})--(\ref{eq:hourly_cap}). The deterministic equivalent is
\begin{equation}
\begin{aligned}
\min_{\bm y,p^{\mathrm{pk}},\{\bm x_s\}}\quad
&C^{\mathrm{ann}}(\bm y,p^{\mathrm{pk}},\{\bm x_s\})\\
\mathrm{s.t.}\quad
&\bm y\in\mathcal Y,\\
&\bm x_s\in\mathcal X(\bm y,\bm\xi_s),\quad s\in\mathcal S.
\end{aligned}
\label{eq:deterministic_equivalent}
\end{equation}
Equation~(\ref{eq:deterministic_equivalent}) is the deterministic equivalent of
CASP and serves as the planning oracle throughout the learning procedure.

% =====================================================================
\section{Decision-Focused Learning Methodology}
\label{sec:methodology}

\subsection{Learning Task and Scenario Generator}

Consider a historical pool $\mathcal D^{\mathrm{tr}}$ whose direct use in every
planning solve is computationally prohibitive. For a support budget
$K\ll|\mathcal D^{\mathrm{tr}}|$, define
$\mathcal I=\{1,\ldots,K\}$. The learning task is to generate $K$ trajectories
that retain the planning information in the full pool rather than merely its
average statistical shape. The conditional decoder produces
\begin{equation}
\widehat{\mathcal S}_\theta=
\{\widehat{\bm\xi}_n=G_\theta(\bm z_n,\bm\omega_n):n\in\mathcal I\},
\label{eq:generated_support}
\end{equation}
where $\bm z_n$ is a latent draw and $\bm\omega_n$ contains cyclical
day-of-year coordinates, weekend share, and mean temperature. The contexts,
latent draws, outage indicators, and support weights are held fixed during
decision-focused fine-tuning. Consequently, a change in the induced design can
be attributed to an update of $G_\theta$, not to resampling. Support weights
preserve the importance-corrected probability mass of normal and outage cases.

The CVAE encoder is
\begin{equation}
q_\phi(\bm z\mid\bm x,\bm\omega)=
\mathcal N(\bm\mu_\phi,\operatorname{diag}(\bm\sigma_\phi^2)),
\label{eq:cvae}
\end{equation}
and the statistical training loss is
\begin{align}
L_{\mathrm{stat}}={}&L_{\mathrm{rec}}+w_{\mathrm{ramp}}L_{\mathrm{ramp}}
+w_{\mathrm{net}}L_{\mathrm{net}}+w_{\mathrm{pk}}L_{\mathrm{pk}}
\nonumber\\
&+w_{\lambda}L_{\lambda\mathrm{-spread}}
+w_{\chi}L_{\chi\mathrm{-spread}}
+\beta L_{\mathrm{KL}}.
\label{eq:cvae_loss}
\end{align}
Here, $L_{\mathrm{rec}}$ is a field-weighted reconstruction loss;
$L_{\mathrm{ramp}}$ preserves hourly changes; and the remaining shape terms
preserve net load, peak net load, price spread, and carbon-intensity spread.
Only $D$, $P^{\mathrm{PV,av}}$, and $\chi$ carry gradients from the decision loss.
Workload and the calendar tariff are restored from admissible templates, while
$A$ is copied as a detached exogenous contingency. The generator therefore
cannot create arbitrage prices or alter outage occurrence to improve its loss.

\subsection{Forward Planning and Fixed-Design Recourse}

Solving the planning oracle on the generated support yields
\begin{equation}
\widehat{\bm y}_\theta\in
\argmin_{\bm y\in\mathcal Y}
C^{\mathrm{ann}}(\bm y;\widehat{\mathcal S}_\theta).
\label{eq:forward_plan}
\end{equation}
Only the first-stage design $\widehat{\bm y}_\theta$ is transferred to the
matched real trajectory $\bm\xi_n$. Its operating variables are discarded and
recourse is solved again:
\begin{equation}
\bm x_n^{F}\in\argmin_{\bm x\in
\mathcal X(\widehat{\bm y}_\theta,\bm\xi_n)}
C^{\mathrm{rec}}(\bm x;\widehat{\bm y}_\theta,\bm\xi_n).
\label{eq:fixed_recourse}
\end{equation}
All dispatch, mode, inventory, curtailment, shedding, and carbon variables are
therefore adapted to the realized trajectory. Let
$\delta^{\mathrm{sh},F}_{n,t}$ and $\delta^{\mathrm{co2},F}_{n,t}$ denote the hourly load
shedding and carbon excess returned by (\ref{eq:fixed_recourse}), and let
$H^F_{n,t,\tau}$ denote the corresponding exposure of delivered carbon in hour
$t$ to grid carbon intensity in hour $\tau$. Solver decisions, violation
severities, and exposures are detached before the loss is formed. They weight
decision-critical prediction errors but do not define a derivative through the
mixed-integer optimizer.

\subsection{Infeasibility-Penalizing Loss}

The infeasibility-penalizing loss (IPL) corrects optimistic generated
constraints only when the induced design fails under real recourse. Define the
normalization terms
\begin{equation}
b^D_{n,t}=\max(|D_{n,t}|,\epsilon_D),\qquad
b^{\mathrm{srv}}_{n,t}=\max(D_{n,t}-\delta^{\mathrm{sh},F}_{n,t},\epsilon_D),
\label{eq:loss_scales}
\end{equation}
and the directional residuals for demand and PV:
\begin{equation}
r^D_{n,t}=\frac{\bigl[D_{n,t}-\widehat D_{n,t}\bigr]_+}{b^D_{n,t}},
\qquad
r^{\mathrm{PV}}_{n,t}=\frac{\bigl[\widehat P^{\mathrm{PV,av}}_{n,t}
-P^{\mathrm{PV,av}}_{n,t}\bigr]_+}
{b^D_{n,t}}.
\label{eq:load_pv_direction}
\end{equation}
These residuals are positive when generated demand is too low or generated PV
availability is too high.

Carbon errors require an intertemporal direction because grid energy charged in
hour $\tau$ may be discharged in hour $t$. To preserve a gradient near the
normal-operation carbon limit, define the smooth carbon pressure
\begin{equation}
\varphi(\chi)=s_\chi\log\left[1+
\exp\left(\frac{\chi-\bar\chi^{\mathrm n}}{s_\chi}\right)\right],
\qquad s_\chi>0.
\label{eq:pressure}
\end{equation}
The exposure-weighted optimistic carbon residual is
\begin{equation}
r^\chi_{n,t}=
\frac{\left[\sum_{\tau\in\mathcal T}H^F_{n,t,\tau}
\bigl[\varphi(\chi_{n,\tau})-\varphi(\widehat\chi_{n,\tau})\bigr]\right]_+}
{\bar\chi^{\mathrm n}\max(\sum_{\tau\in\mathcal T}H^F_{n,t,\tau},\epsilon_H)}.
\label{eq:carbon_direction}
\end{equation}
For the zero-anchored smooth margin
$h_m(r)=\operatorname{softplus}(m+r)-\operatorname{softplus}(m)$, the three IPL
components are
\begin{align}
L_D^{\mathrm{IPL}}&=\frac1{KT}\sum_{n\in\mathcal I}
\sum_{t\in\mathcal T}
\frac{\delta^{\mathrm{sh},F}_{n,t}}{b^D_{n,t}}h_m(r^D_{n,t}),
\label{eq:ipl_load}\\
L_{\mathrm{PV}}^{\mathrm{IPL}}&=\frac1{KT}\sum_{n\in\mathcal I}
\sum_{t\in\mathcal T}
\frac{\delta^{\mathrm{sh},F}_{n,t}}{b^D_{n,t}}h_m(r^{\mathrm{PV}}_{n,t}),
\label{eq:ipl_pv}\\
L_\chi^{\mathrm{IPL}}&=\frac1{KT}\sum_{n\in\mathcal I}
\sum_{t\in\mathcal T}
\frac{\delta^{\mathrm{co2},F}_{n,t}}{b^{\mathrm{srv}}_{n,t}}h_m(r^\chi_{n,t}),
\label{eq:ipl_carbon}
\end{align}
and
\begin{equation}
L_{\mathrm{IPL}}=w_D L_D^{\mathrm{IPL}}
+w_{\mathrm{PV}}L_{\mathrm{PV}}^{\mathrm{IPL}}
+w_\chi L_\chi^{\mathrm{IPL}}.
\label{eq:ipl}
\end{equation}
Gradient descent therefore raises underestimated demand, lowers overestimated
PV availability, and raises underestimated carbon pressure only at locations
weighted by real recourse violations.

\subsection{Optimality-Preserving Loss}

IPL alone can reduce violations by making every generated scenario severe. To
prevent this degenerate solution, let
$\mathcal B^R=\{\bm\xi_n:n\in\mathcal I\}$ be the matched real batch and solve
\begin{equation}
\bm y^\star\in\argmin_{\bm y\in\mathcal Y}
C^{\mathrm{ann}}(\bm y;\mathcal B^R).
\label{eq:pi_design}
\end{equation}
The reference design $\bm y^\star$ is fixed and re-dispatched on
$\widehat{\mathcal S}_\theta$. Denote the resulting shedding, carbon excess,
and carbon-exposure matrix by $\delta^{\mathrm{sh},P}_{n,t}$,
$\delta^{\mathrm{co2},P}_{n,t}$, and $H^P_{n,t,\tau}$, respectively. A positive violation
means that the generated constraints have excluded the design optimized for the
matched real batch.

The OPL uses the reverse directional residuals
\begin{equation}
\widetilde r^D_{n,t}=
\frac{\bigl[\widehat D_{n,t}-D_{n,t}\bigr]_+}{b^D_{n,t}},
\qquad
\widetilde r^{\mathrm{PV}}_{n,t}=
\frac{\bigl[P^{\mathrm{PV,av}}_{n,t}
-\widehat P^{\mathrm{PV,av}}_{n,t}\bigr]_+}
{b^D_{n,t}},
\label{eq:opl_directions}
\end{equation}
and the exposure-weighted carbon residual
\begin{equation}
\widetilde r^\chi_{n,t}=
\frac{\left[\sum_{\tau\in\mathcal T}H^P_{n,t,\tau}
\bigl[\varphi(\widehat\chi_{n,\tau})-\varphi(\chi_{n,\tau})\bigr]\right]_+}
{\bar\chi^{\mathrm n}\max(\sum_{\tau\in\mathcal T}H^P_{n,t,\tau},\epsilon_H)}.
\label{eq:opl_carbon_direction}
\end{equation}
With $\widetilde b^{\mathrm{srv}}_{n,t}=\max(\widehat D_{n,t}
-\delta^{\mathrm{sh},P}_{n,t},\epsilon_D)$, define
\begin{align}
L_D^{\mathrm{OPL}}&=\frac1{KT}\sum_{n\in\mathcal I}
\sum_{t\in\mathcal T}
\frac{\delta^{\mathrm{sh},P}_{n,t}}{b^D_{n,t}}
h_m(\widetilde r^D_{n,t}),
\label{eq:opl_load}\\
L_{\mathrm{PV}}^{\mathrm{OPL}}&=\frac1{KT}\sum_{n\in\mathcal I}
\sum_{t\in\mathcal T}
\frac{\delta^{\mathrm{sh},P}_{n,t}}{b^D_{n,t}}
h_m(\widetilde r^{\mathrm{PV}}_{n,t}),
\label{eq:opl_pv}\\
L_\chi^{\mathrm{OPL}}&=\frac1{KT}\sum_{n\in\mathcal I}
\sum_{t\in\mathcal T}
\frac{\delta^{\mathrm{co2},P}_{n,t}}{\widetilde b^{\mathrm{srv}}_{n,t}}
h_m(\widetilde r^\chi_{n,t}).
\label{eq:opl_carbon}
\end{align}
The complete preservation loss is
\begin{equation}
L_{\mathrm{OPL}}=w_D L_D^{\mathrm{OPL}}
+w_{\mathrm{PV}}L_{\mathrm{PV}}^{\mathrm{OPL}}
+w_\chi L_\chi^{\mathrm{OPL}}.
\label{eq:opl}
\end{equation}
It lowers excessively high demand and carbon predictions and raises
excessively low PV availability only where the generated support rejects the
reference design.

Finally, the two directional losses are combined as
\begin{equation}
L_{\mathrm{dec}}=\alpha L_{\mathrm{IPL}}+(1-\alpha)L_{\mathrm{OPL}},
\qquad 0\le\alpha\le1.
\label{eq:dfl_loss}
\end{equation}
This construction adapts the feasibility/optimality logic of
~\cite{Mandi2025} to two-stage fixed-design recourse. It is a directional
surrogate, not a derivative of the MILP value function.

\subsection{Gradient Balancing and Training}

Decision feedback and statistical reconstruction are evaluated on independent
batches. Their combined objective is
\begin{equation}
L_{\mathrm{total}}=w_{\mathrm{stat}}L_{\mathrm{stat}}
+w^{\mathrm{eff}}_{\mathrm{dec}}L_{\mathrm{dec}}.
\label{eq:total_loss}
\end{equation}
Let $g_{\mathrm{stat}}$ and $g_{\mathrm{dec}}$ be the corresponding decoder
gradient norms. The detached decision-loss weight is
\begin{equation}
w^{\mathrm{eff}}_{\mathrm{dec}}=w_{\mathrm{dec}}
\operatorname{clip}\left(
r_g\frac{g_{\mathrm{stat}}}{\max(g_{\mathrm{dec}},\epsilon_g)},
b_{\min},b_{\max}\right).
\label{eq:gradient_balance}
\end{equation}
Only decoder parameters are used in this comparison because both loss terms
reach that parameter subspace, whereas only the statistical loss reaches the
encoder.

\begin{algorithm}[H]
\caption{Feasibility-aware CVAE fine-tuning}
\label{alg:training}
\begin{algorithmic}[1]
\Require CVAE, training and validation pools, planning oracle, support size $K$
\State Fix support contexts, latent draws, outage indicators, and weights
\For{each fine-tuning epoch}
  \State Decode $\widehat{\mathcal S}_\theta$ and solve (\ref{eq:forward_plan})
  \State Re-dispatch $\widehat{\bm y}_\theta$ on $\mathcal B^R$
  \State Solve (\ref{eq:pi_design}) and re-dispatch $\bm y^\star$ on
  $\widehat{\mathcal S}_\theta$
  \State Form $L_{\mathrm{IPL}}$, $L_{\mathrm{OPL}}$, and $L_{\mathrm{dec}}$
  \State Evaluate $L_{\mathrm{stat}}$ on an independent reconstruction batch
  \State Balance decoder gradients and update $\theta$
  \State Evaluate the fixed support on the validation pool
\EndFor
\State Restore the best validation checkpoint and return its $K$ scenarios
\end{algorithmic}
\end{algorithm}

Checkpoint selection is lexicographic. Checkpoints that exceed the prescribed
shedding tolerance are rejected; among reliable checkpoints, lower carbon
excess is preferred and exact decision regret breaks ties. The initial model is
retained as an epoch-$-1$ candidate, so decision-focused updates are accepted
only when they improve the fixed validation criterion.

\subsection{Out-of-Sample Metrics}

Let $\mathcal S^{\mathrm{te}}$ be the common held-out test set. Its
perfect-information cost and the fixed-design cost of method $m$ are
\begin{align}
J^{\mathrm{PI}}&=\min_{\bm y\in\mathcal Y}
C^{\mathrm{ann}}(\bm y;\mathcal S^{\mathrm{te}}),
\label{eq:test_pi}\\
J_m&=C^{\mathrm{ann}}(\bm y_m;\mathcal S^{\mathrm{te}}),
\label{eq:test_cost}\\
\mathcal R_m&=\frac{J_m-J^{\mathrm{PI}}}{|J^{\mathrm{PI}}|+\epsilon_J}.
\label{eq:regret}
\end{align}
The regret $\mathcal R_m$ is an evaluation metric and never enters gradient
calculation. With $J^0$ denoting the no-storage cost on the same test set, the
storage value is $\mathcal V_m=J^0-J_m$. Annual cost is reported together with
load shedding, carbon excess, PV curtailment, and the installed power and energy
ratings.

% =====================================================================
\section{Case Study}
\label{sec:case}

\subsection{Data and Experimental Setup}

The dataset is divided into nonoverlapping 48-hour trajectories covering three
years. Facility demand is calibrated to a 10-MW-class data center using a
workload proxy and temperature-dependent power-usage effectiveness. The PV
shape is retained from the source feeder location and rescaled to a 5-MW
installation. Electricity tariffs follow their calendar templates, and each
grid-carbon trajectory remains an intact 48-hour block. Any within-season
permutation is performed separately inside the training, validation, and test
splits, which remain mutually disjoint.

Contiguous outages of one to four hours are injected into a stratified subset
of trajectories so that small decision batches contain reliability events.
Because outages are deliberately oversampled for learning, importance weights
restore an assumed physical probability of approximately $0.71\%$ per 48-hour
window. Outage time is sampled independently of demand, PV, and carbon; it is
therefore a controlled stress-test condition rather than an estimate of their
joint field distribution.

\begin{table}[t]
\centering
\caption{Principal system and planning parameters}
\label{tab:parameters}
\begin{tabular}{@{}lr@{}}
\toprule
Parameter & Value \\
\midrule
Scenario horizon/resolution & 48 h/1 h \\
Mean/maximum demand & 8.477/10.224 MW \\
PV capacity & 5 MW \\
Grid-import limit & 15 MW \\
Backup-generator capacity & 6 MW \\
Battery power bound if built & 0.2--7.5 MW \\
Battery energy bound if built & 0.5--60 MWh \\
Initial/minimum/maximum SOC & 0.50/0.10/0.90 \\
Charge/discharge efficiency & 0.95/0.95 \\
Normal/outage carbon limit & 0.28/0.66 tCO$_2$/MWh \\
Planning support size & 8 \\
\bottomrule
\end{tabular}
\end{table}

The CVAE has a 16-dimensional latent space and two 256-unit hidden layers in
both encoder and decoder. Statistical pretraining runs for 300 epochs and
decision-focused fine-tuning for 45 epochs. The principal configuration fixes
eight decision anchors, uses an independent 32-trajectory reconstruction batch,
sets $\alpha=0.5$, and targets a decoder-gradient ratio of $r_g=0.1$. Planning
solves used during training allow a $1\%$ relative gap; final evaluation uses
the tighter oracle tolerance.

\subsection{Compared Methods and Metrics}

The statistical-selection baselines are random sampling, K-means,
farthest-point selection, and an aggregate observed support. Learning baselines
include a statistically trained CVAE, IPL-only fine-tuning, and OPL-only
fine-tuning; the proposed method uses both directional losses. Every method
receives the same support budget and is evaluated by fixing its selected design
and solving recourse on the identical test trajectories. Direct optimization on
the full test set is used only to define the perfect-information reference.

\begin{table*}[t]
\centering
\caption{Out-of-sample planning performance on a common test set}
\label{tab:main_results}
\begin{tabular}{@{}lrrrrrrr@{}}
\toprule
Method & $P^{\mathrm{ES}}$ & $E^{\mathrm{ES}}$ & Annual cost & $\mathcal R$
& Shedding & Carbon excess & Storage value \\
& (MW) & (MWh) & (\$) & (\%) & (MWh) & (tCO$_2$) & (\$) \\
\midrule
Random & -- & -- & -- & -- & -- & -- & -- \\
K-means & -- & -- & -- & -- & -- & -- & -- \\
Farthest & -- & -- & -- & -- & -- & -- & -- \\
Aggregate & -- & -- & -- & -- & -- & -- & -- \\
CVAE-only & -- & -- & -- & -- & -- & -- & -- \\
IPL-only & -- & -- & -- & -- & -- & -- & -- \\
OPL-only & -- & -- & -- & -- & -- & -- & -- \\
Full DFL & -- & -- & -- & -- & -- & -- & -- \\
Perfect information & -- & -- & -- & 0 & -- & -- & -- \\
\bottomrule
\end{tabular}
\end{table*}

% TODO: Populate from multi-seed evaluation artifacts. Report mean and standard
% deviation and distinguish incumbents from certified bounds where necessary.

The completed analysis will first compare the selected power and energy ratings,
then determine whether cost differences are accompanied by load shedding or
carbon excess, and finally compare the observed differences with both
multi-seed variability and the uncertainty implied by solver gaps.

\subsection{Carbon-Accounting and Learning Ablations}

Four matched planning models compare source-resolved and perfectly mixed carbon
accounting under hourly and horizon-average carbon limits. Data, scenario
weights, cost coefficients, design bounds, and solver settings are held fixed
across all four models.

\begin{table}[t]
\centering
\caption{Effect of carbon accounting and cap scope}
\label{tab:carbon_ablation}
\begin{tabular}{@{}llrrr@{}}
\toprule
Accounting & Scope & $P^{\mathrm{ES}}$ & $E^{\mathrm{ES}}$ & Excess \\
& & (MW) & (MWh) & (tCO$_2$) \\
\midrule
Layered & Hourly & -- & -- & -- \\
Layered & Horizon & -- & -- & -- \\
Average & Hourly & -- & -- & -- \\
Average & Horizon & -- & -- & -- \\
\bottomrule
\end{tabular}
\end{table}

% TODO: Current t1--t4 YAML files are cost/DFL sensitivities, not the four
% accounting-scope combinations. Create matched configurations first.

The learning ablation varies $\alpha$, removes decoder-gradient balancing, and
changes the support size. Planning sensitivities vary the normal-operation
carbon limit, battery capital cost, and outage duration. Multi-seed results are
reported as means and standard deviations. A representative 48-hour dispatch
will show grid and delivered carbon intensity, aggregate state of charge,
charging source, and discharge by carbon layer so that differences in selected
capacity can be traced to operating behavior.

\subsection{Computational Performance and Limitations}

Computational reporting separates forward-planning, fixed-design recourse, and
matched-real reference solves; it also records solver calls per epoch, wall
time, and final optimality gaps. The aggregate facility model cannot support
claims about network congestion or multi-location siting. Demand is calibrated
rather than metered, outage time is simulated independently of weather, and the
directional decision loss is not a derivative of the MILP value function.

% =====================================================================
\section{Conclusion}
\label{sec:conclusion}

This paper has developed a two-stage formulation for carbon-constrained battery
planning in a grid-connected data center. The model jointly accounts for
investment, energy purchases, demand charges, storage degradation, outage-only
load shedding, and the carbon intensity delivered to the facility. Virtual
inventory layers retain both energy source and grid-charging time while sharing
the physical power and energy limits of one battery. The resulting ledger is
linear and avoids the bilinear intensity identities of a perfectly mixed carbon
inventory.

A feasibility-aware learning method has also been developed to construct a
compact planning support. In the forward direction, the design selected from
generated scenarios is fixed and re-dispatched on matched real trajectories; in
the preservation direction, a matched-real reference design is fixed and
re-dispatched on generated trajectories. The corresponding shedding, carbon
excess, and carbon-exposure outputs distinguish optimism from excessive
conservatism. Because these solver outputs are detached, the method retains
gradients to generated demand, PV availability, and carbon intensity without
differentiating through the mixed-integer oracle.

% TODO: Insert quantitative findings after the result tables are populated.

Future work will consider network-constrained multi-location planning,
weather-correlated outages, flexible computing demand, and longer horizons with
linked inter-block storage states.

\section*{Acknowledgements}
% Add funding and non-author contributions.

\section*{CRediT authorship contribution statement}
Author Name: Conceptualization, Methodology, Software, Validation,
Writing--original draft.

\section*{Declaration of competing interest}
The authors declare that they have no known competing financial interests or
personal relationships that could have appeared to influence the work reported
in this paper.

\section*{Data availability}
The processed data, configurations, and source code used to reproduce the study
will be made available with the published article.

\bibliographystyle{elsarticle-num}
\bibliography{main_elsevier}
\end{document}
